% ============================================================
% M3 S2 导学件·波1 片件 —— 课时06B 点到直线距离与平行线距离（2.2.4·选学补位）
% 生成：M3 成卷轮 S2 波1 臂3｜2026-09-14｜编译 cwd＝本目录（中文路径禁 \input 绝对路径）
%   xelatex main-true.tex ×2 → 含详解印本；xelatex main-false.tex ×2 → 纯题版（单源双本）
% 输入链（全只读）：题面库 课时06B-2.2.4点到直线距离.md（题面逐字装配）＋ -答案侧.md（值逐字回嵌）
%   ＋ 定稿/批B补-课时06B-2.2.4 点到直线距离.md（详解回嵌源）；toolchain＝qp-m3.sty（钉后版
%   md5 7c3930362be8a0a2bdf21bbf8ac16573，本地挂载副本，破损即停）。
% 母版体例：照抄 成卷/导学件/课时01-坐标法/（骨架/宏用法/符号路由 preamble 整块照抄）。
% 片级注记承接（批B 台账下发）：①本件系选学补位课时（2.2.4），序位于课时06 与课时07 之间，
%   件头标选学注记（批B-课时06 T13/T16 前引注同口径）；②批B 台账 §五.4 勘误
%   「8x→8y」落点在本片下游引用口径（按「8y」取，本片题面无涉）；③T15 导数符号分支
%   改写系课时07 注记，本片无涉。
% 键账：件内 ansblock 键集 ≡ 题面库片键集 ≡ manifest 键序（21 键，零缺漏零多余）；
%   预习位零键零答案（口令④裁定前口径）；印面号 1—21＝探究 1—16＋课堂评价 17—21，
%   键↔号映射见 值台账-课时06B.json。
% 括线判据（承重墙禁放宽）：估高 top-5（wlen 全角1·ASCII0.5，剔\命令／23 栏宽字口径）
%   → 括线模，逐键读数入值台账；其余灰底。本片括线键见 值台账-课时06B.json items 判模列。
% 门谱读数：见 过程件 工作区/_tmpM3S2波1臂3_0914/（编译 log、门谱输出、目验 PNG 双档）。
% ============================================================
\documentclass[fontset=none]{ctexart}
\usepackage{qp-m3}
% —— 印面逐字符号路由补钉（照抄课时01 母版 0914 实证块；禁改 sty 本体保 md5） ——
\xeCJKDeclareCharClass{Default}{"2212}%
\setCJKfamilyfont{nscblk}[Path=C:/提示词/工作区/字替对照-0909/variantF/fonts/,BoldFont=NSC-w600.ttf]{NSC-w500.ttf}%
\xeCJKDeclareSubCJKBlock{nscblk}{"25B1}%
\newcommand{\pxparallelogram}{{\CJKfamily{nscblk}▱}}%
% —— 断行微调（照抄母版：CJK 胶收缩 0.01→0.05em；\emergencystretch 不设） ——
\xeCJKsetup{CJKglue={\hskip 0.10em plus 0.08em minus 0.05em}}%
% —— 页脚件名（qp-headfoot 复用入口） ——
\renewcommand{\qpzhangming}{第二章\quad 平面解析几何}%
\renewcommand{\qpjianming}{导学件}%

\begin{document}
%装配轮S6三本跨本连续页码（G7方案B）setcounter 注入位
\setcounter{page}{21}

\zhangtitle{第二章\quad 平面解析几何}
\jietitle{2.2.4 点到直线距离与平行线距离（选学）}
\zhuzhu{选学注记：本课时节为选学补位（人教B版§2.2.4），序位在课时06（§2.2.3）与课时07（§2.3.1）之间，可与§2.2.3 交点坐标配套使用；课时06 件内 T13/T16 已加「需选学§2.2.4(06B)」前引注.}
\mubiaoline
\mubiaomu{1}{理解点到直线距离公式的推导思想，掌握点到直线距离公式与两平行直线间距离公式，并能正确应用；}
\mubiaomu{2}{会用距离公式求参数、求三角形面积，能把\(m^{2}+n^{2}\)、\(\sqrt{(m-a)^{2}+(n-b)^{2}}\)等代数式的最值转化为距离问题；}
\mubiaomu{3}{掌握点（线）关于直线的对称与光线反射模型，体会“对称变换、几何最短”的转化思想.}

\vspace{1.7mm}
\begin{multicols}{2}
\emergencystretch=1em

% ============================================================
% 课前预习（知识点导学填空；零键零答案——预习键位按检查口令④裁定前口径，
% \kongbai 空线＋\zhentib 判断题面形，两档等形不泄答）
% ============================================================
\huaxing{课}{前}{预}{习}{知识导学\quad 素养初识}

\zsd{一}{点到直线的距离}

\tiaomu{1}{点\(P(x_0,y_0)\)到直线\(l\)：\(ax+by+c=0\)（\(a^{2}+b^{2}\neq 0\)）的距离\(d=\)\kongbai{}；使用前须把直线方程化为\kongbai{}式，并核对系数．\zhuzhu{距离公式对含参直线同样适用：先按公式表示距离，再按距离条件列方程（组）反求参数，解出后回代检验（排除重合、退化情形）.}}

\zsd{二}{两平行直线间的距离}

\tiaomu{1}{两平行直线\(Ax+By+C_1=0\)与\(Ax+By+C_2=0\)（\(x\)，\(y\)系数已化同）间的距离\(d=\)\kongbai{}．\zhuzhu{系数不同先化同再套公式；平行线间的距离是两线间的公共垂线段长.}}

\tiaomu{2}{与两平行直线等距离（或距离成定比）的直线，常设平行线系\(Ax+By+m=0\)，由距离条件解出\(m\)，注意分侧讨论、逐解检验．\zhuzhu{距离成比的问题带域内、带域外都要讨论，漏解常丢在内侧一档.}}

\zsd{三}{对称与光线反射}

\tiaomu{1}{点\(P(x_0,y_0)\)关于直线\(l\)的对称点\(P'\)由两条约束确定：\(PP'\)\kongbai{}\(l\)，且\(PP'\)的中点\kongbai{}；光线反射求最短路程，把入射点换成定点关于反射面的对称点，化为两点间距离．\zhuzhu{“垂直＋中点在对称轴上”是对称点的通用方程组；反射问题（入射角＝反射角）与对称最短是同一模型.}}

\zhenhead{判断正误(正确的打\gou{},错误的打\cha{})}

\zhentib{(1)}{点\(P\)到直线\(l\)的距离就是点\(P\)与直线\(l\)上任意一点连线的长度．}

\zhentib{(2)}{两条平行直线间的距离等于两方程常数项之差的绝对值除以一次项系数平方和的算术平方根，与系数是否化同无关．}

\zhentib{(3)}{光线经直线反射，入射光线所在直线与反射光线所在直线关于反射面对称．}

\liubai[9mm]{此处书写}

% ============================================================
% 课中探究（考点探究〔例+变式〕＝练习槽 16 键；题后紧跟答案＝ansblock 内嵌，
% 值照答案侧逐字，详解承批B补 亲算回嵌＋印面清洗；装配序＝印面号 1—16）
% ============================================================
\huaxing{课}{中}{探}{究}{考点探究\quad 素养小结}

% ---------- 探究点一 距离公式直用与几何意义 ----------
\tjdnr{一}{距离公式直用与几何意义}{例\textbf{1}}{中档(距离方程两解)}{求在\(y\)轴上且到直线\(3x+4y+5=0\)的距离等于5的点的坐标\nobreak (\kongwei)}

\bindopt {\fontsize{10.5pt}{\qpTJgLeadLiti}\selectfont \makebox[\dimexpr0.5\linewidth-1em\relax][l]{A．\((0,5)\)}\makebox[\dimexpr0.5\linewidth-1em\relax][l]{B．\((0,-15/2)\)}\par}

\bindopt {\fontsize{10.5pt}{\qpTJgLeadLiti}\selectfont \makebox[\dimexpr0.5\linewidth-1em\relax][l]{C．\((0,5)\)或\((0,-15/2)\)}\makebox[\dimexpr0.5\linewidth-1em\relax][l]{D．\((5,0)\)或\((-5,0)\)}\par}

\begin{ansblock}[2章-练-课时06B-T1]
% ans:2章-练-课时06B-T1
\ansitem{1}{C}
\ansnote{详解}{设点\((0,t)\)，\(d=|4t+5|/\sqrt{3^{2}+4^{2}}=|4t+5|/5=5\)，即\(|4t+5|=25\)，解得\(t=5\)或\(t=-15/2\)。所求点为\((0,5)\)或\((0,-15/2)\)，故选：C．（D 中\((5,0)\)、\((-5,0)\)到直线的距离分别为4与2，不合，系轴别混淆干扰项．）}
\end{ansblock}

\liB{变式\textbf{2}}{中档(三角形面积)}{}{已知点\(A(1,3)\)，\(B(2,1)\)，\(C(-1,0)\)，则\(\triangle ABC\)的面积为\nobreak (\kongwei)}

\bindopt {\fontsize{10.5pt}{\qpTJgLeadLiti}\selectfont \makebox[\dimexpr0.25\linewidth-0.5em\relax][l]{A．\(7\)}\makebox[\dimexpr0.25\linewidth-0.5em\relax][l]{B．\(13/2\)}\makebox[\dimexpr0.25\linewidth-0.5em\relax][l]{C．\(5\)}\makebox[\dimexpr0.25\linewidth-0.5em\relax][l]{D．\(7/2\)}\par}

\begin{ansblock}[2章-练-课时06B-T3]
% ans:2章-练-课时06B-T3
\ansitem{2}{D}
\ansnote{详解}{以\(BC\)为底：\(|BC|=\sqrt{(-1-2)^{2}+(0-1)^{2}}=\sqrt{10}\)；直线\(BC\)过\(B(2,1)\)、\(C(-1,0)\)，斜率为\(1/3\)，方程\(x-3y+1=0\)。\(A(1,3)\)到\(BC\)的高\(h=|1-9+1|/\sqrt{1^{2}+(-3)^{2}}=7/\sqrt{10}\)。故\(S=(1/2)\times\sqrt{10}\times 7/\sqrt{10}=7/2\)，选：D．}
\end{ansblock}

\liB{变式\textbf{3}}{中档(点斜式＋距离公式)}{}{若直线\(l\)过点\((2,\ \sqrt{3})\)，倾斜角为\(120^\circ\)，则点\((1,\ -\sqrt{3})\)到直线\(l\)的距离为\nobreak (\kongwei)}

\bindopt {\fontsize{10.5pt}{\qpTJgLeadLiti}\selectfont \makebox[\dimexpr0.5\linewidth-1em\relax][l]{A．\(\sqrt{3}/2\)}\makebox[\dimexpr0.5\linewidth-1em\relax][l]{B．\(\sqrt{3}\)}\par}

\bindopt {\fontsize{10.5pt}{\qpTJgLeadLiti}\selectfont \makebox[\dimexpr0.5\linewidth-1em\relax][l]{C．\(3\sqrt{3}/2\)}\makebox[\dimexpr0.5\linewidth-1em\relax][l]{D．\(5\sqrt{3}/2\)}\par}

\begin{ansblock}[2章-练-课时06B-T9]
% ans:2章-练-课时06B-T9
\ansitem{3}{C}
\ansnote{详解}{\(k=\tan 120^\circ=-\sqrt{3}\)，直线\(l\)：\(y-\sqrt{3}=-\sqrt{3}(x-2)\)，即\(\sqrt{3}x+y-3\sqrt{3}=0\)。\(d=|\sqrt{3}\times 1-\sqrt{3}-3\sqrt{3}|/\sqrt{(\sqrt{3})^{2}+1^{2}}=3\sqrt{3}/2\)，故选：C．}
\end{ansblock}

\liB{变式\textbf{4}}{中档(垂线段最小·几何意义)}{}{若点\(M(m,n)\)为直线\(l\)：\(3x+4y+2=0\)上的动点，则\(m^{2}+n^{2}\)的最小值为\kongbai{}．}
\xiexwei{7mm}

\begin{ansblock}[2章-练-课时06B-T12]
% ans:2章-练-课时06B-T12
\ansitem{4}{4/25}
\ansnote{详解}{\(m^{2}+n^{2}=|OM|^{2}\)，其最小值在\(OM\)垂直于\(l\)时取得，即原点到\(l\)的距离：\(d=|0+0+2|/\sqrt{3^{2}+4^{2}}=2/5\)。故\(m^{2}+n^{2}\)的最小值为\(d^{2}=4/25\)．}
\end{ansblock}

\tjdnr{一}{距离公式直用与几何意义}{例\textbf{5}}{中档(含参距离最大·辅助角)}{点\(A(1,1)\)到直线\(x\cos\theta+y\sin\theta-2=0\)的距离的最大值是\kongbai{}．}
\xiexwei{7mm}

\begin{ansblock}[2章-练-课时06B-T13]
% ans:2章-练-课时06B-T13
\ansitem{5}{2+√2}
\ansnote{详解}{\(d=|\cos\theta+\sin\theta-2|/\sqrt{\cos^{2}\theta+\sin^{2}\theta}\)\par\noindent \(=|\sqrt{2}\sin(\theta+\pi/4)-2|\)。\par\noindent 因\(\sqrt{2}\sin(\theta+\pi/4)\leqslant\sqrt{2}<2\)，绝对值内恒为负。\par\noindent 故\(d=2-\sqrt{2}\sin(\theta+\pi/4)\)，当\(\sin(\theta+\pi/4)=-1\)时取最大值\(2+\sqrt{2}\)．}
\end{ansblock}

\xiaojie{①识别：以求距离、由距离反求参数、或把代数式最值翻译成距离为目标的题，判归本题型；含参直线先表示距离再列方程.}
\bindp {\kaishu ②操作：套公式前先化一般式并核对系数；平方去绝对值须非负前提＋逐根回代检验；\(m^{2}+n^{2}\)型看原点距离，\(\sqrt{(m-a)^{2}+(n-b)^{2}}\)型看两点距离.}
\bindp {\kaishu ③收束：含参三角距离用辅助角收拢成\(R\sin(\theta+\varphi)\)型，先判绝对值内符号再去绝对值，最值在端点或峰值处取得.}

% ---------- 探究点二 平行线间的距离与最值 ----------
\tjdnr{二}{平行线间的距离与最值}{例\textbf{6}}{中档(动平行线距离范围)}{已知两平行直线\(l_1\)，\(l_2\)分别过点\(P(-1,3)\)，\(Q(2,-1)\)，它们分别绕\(P\)，\(Q\)旋转，但始终保持平行，则\(l_1\)，\(l_2\)之间的距离的取值范围是\nobreak (\kongwei)}

\bindopt {\fontsize{10.5pt}{\qpTJgLeadLiti}\selectfont \makebox[\dimexpr0.5\linewidth-1em\relax][l]{A．\((0,+\infty)\)}\makebox[\dimexpr0.5\linewidth-1em\relax][l]{B．\([0,5]\)}\par}

\bindopt {\fontsize{10.5pt}{\qpTJgLeadLiti}\selectfont \makebox[\dimexpr0.5\linewidth-1em\relax][l]{C．\((0,5]\)}\makebox[\dimexpr0.5\linewidth-1em\relax][l]{D．\([0,\sqrt{17}]\)}\par}

\begin{ansblock}[2章-练-课时06B-T5]
% ans:2章-练-课时06B-T5
\ansitem{6}{C}
\ansnote{详解}{两线过\(P\)、\(Q\)且始终平行。当两线均与\(PQ\)垂直时距离最大，最大值\(=|PQ|=\sqrt{(2+1)^{2}+(-1-3)^{2}}=5\)；当两线向重合方向旋转时距离趋于0，但重合时不能分别过两不同点，0取不到。故\(d\in(0,5]\)，故选：C．}
\end{ansblock}

\liB{变式\textbf{7}}{中档(韦达与距离最值)}{}{设两条直线的方程分别为\(x+y+a=0\)，\(x+y+b=0\)，已知\(a\)，\(b\)是方程\(x^{2}+x+c=0\)的两个实根，且\(0\leqslant c\leqslant 1/8\)，则这两条直线之间的距离的最大值和最小值分别为\nobreak (\kongwei)}

\bindopt {\fontsize{10.5pt}{\qpTJgLeadLiti}\selectfont \makebox[\dimexpr0.5\linewidth-1em\relax][l]{A．\(\sqrt{3}/3\)，\(\sqrt{2}/3\)}\makebox[\dimexpr0.5\linewidth-1em\relax][l]{B．\(\sqrt{3}/3\)，\(1/3\)}\par}

\bindopt {\fontsize{10.5pt}{\qpTJgLeadLiti}\selectfont \makebox[\dimexpr0.5\linewidth-1em\relax][l]{C．\(\sqrt{2}/2\)，\(1/2\)}\makebox[\dimexpr0.5\linewidth-1em\relax][l]{D．\(\sqrt{2}/2\)，\(\sqrt{2}/3\)}\par}

\begin{ansblock}[2章-练-课时06B-T6]
% ans:2章-练-课时06B-T6
\ansitem{7}{C}
\ansnote{详解}{两线平行，\(d=|a-b|/\sqrt{2}\)。由韦达定理\(a+b=-1\)，\(ab=c\)，故\(|a-b|=\sqrt{(a+b)^{2}-4ab}=\sqrt{1-4c}\)。由\(0\leqslant c\leqslant 1/8\)得\(1/2\leqslant 1-4c\leqslant 1\)，故\(\sqrt{2}/2\leqslant|a-b|\leqslant 1\)。于是\(d\in[1/2,\ \sqrt{2}/2]\)：最大值\(\sqrt{2}/2\)，最小值\(1/2\)，故选：C．（实根存在性：\(\Delta=1-4c\geqslant 0\)由\(c\leqslant 1/8\)保证．）}
\end{ansblock}

\liB{变式\textbf{8}}{中档(平行线间垂线段转化)}{}{已知点\(P\)，\(Q\)分别在直线\(l_1\)：\(x+y+2=0\)与直线\(l_2\)：\(x+y-1=0\)上，且\(PQ\perp l_1\)，点\(A(-3,-3)\)，\(B(3/2,\ 1/2)\)，则\(|AP|+|PQ|+|QB|\)的最小值为\nobreak (\kongwei)}

\bindopt {\fontsize{10.5pt}{\qpTJgLeadLiti}\selectfont \makebox[\dimexpr0.5\linewidth-1em\relax][l]{A．\(\sqrt{130}/2\)}\makebox[\dimexpr0.5\linewidth-1em\relax][l]{B．\(\sqrt{13}+3\sqrt{2}/2\)}\par}

\bindopt {\fontsize{10.5pt}{\qpTJgLeadLiti}\selectfont \makebox[\dimexpr0.5\linewidth-1em\relax][l]{C．\(\sqrt{13}\)}\makebox[\dimexpr0.5\linewidth-1em\relax][l]{D．\(3\sqrt{2}\)}\par}

{\ansblockgrayfalse
\begin{ansblock}[2章-练-课时06B-T8]
% ans:2章-练-课时06B-T8
\ansitem{8}{B}
\ansnote{详解}{两线平行，\(|PQ|\)等于平行线间的距离，为定值\(=|2-(-1)|/\sqrt{2}=3\sqrt{2}/2\)；且向量\(\overrightarrow{QP}\)恒为\(l_1\)指向\(l_2\)的法向位移\((3/2,\ 3/2)\)。作平移\(A'=A+(3/2,\ 3/2)=(-3/2,\ -3/2)\)，则\(AA'\parallel QP\)且\(|AA'|=|QP|\)，四边形\(AA'QP\)为平行四边形，故\(|AP|=|A'Q|\)。于是\(|AP|+|PQ|+|QB|=|A'Q|+3\sqrt{2}/2+|QB|\geqslant|A'B|+3\sqrt{2}/2\)，当\(A'\)、\(Q\)、\(B\)共线（\(Q\)在线段上）时取等。\(|A'B|=\sqrt{(3/2+3/2)^{2}+(1/2+3/2)^{2}}=\sqrt{13}\)；取等点验证：线段\(A'B\)上参数\(t=4/5\)处的点\((9/10,\ 1/10)\)满足\(9/10+1/10-1=0\)，确在\(l_2\)上，等号可取。故最小值\(=\sqrt{13}+3\sqrt{2}/2\)，故选：B．}
\end{ansblock}}

\tjdnr{二}{平行线间的距离与最值}{例\textbf{9}}{简单(平行线距离求参·多解)}{若两条平行直线\(l_1\)：\(x-2y+m=0\)与\(l_2\)：\(2x+ny-6=0\)之间的距离是\(2\sqrt{5}\)，则\(m+n\)的可能值为\nobreak (\kongwei)}

\bindopt {\fontsize{10.5pt}{\qpTJgLeadLiti}\selectfont \makebox[\dimexpr0.25\linewidth-0.5em\relax][l]{A．\(3\)}\makebox[\dimexpr0.25\linewidth-0.5em\relax][l]{B．\(-17\)}\makebox[\dimexpr0.25\linewidth-0.5em\relax][l]{C．\(-3\)}\makebox[\dimexpr0.25\linewidth-0.5em\relax][l]{D．\(17\)}\par}

\begin{ansblock}[2章-练-课时06B-T10]
% ans:2章-练-课时06B-T10
\ansitem{9}{AB}
\ansnote{详解}{平行\(\iff 1/2=(-2)/n\)（\(n\neq 0\)），得\(n=-4\)。\(l_2\)：\(2x-4y-6=0\)，即\(x-2y-3=0\)。\(d=|m-(-3)|/\sqrt{1^{2}+(-2)^{2}}=|m+3|/\sqrt{5}=2\sqrt{5}\)，得\(|m+3|=10\)，\(m=7\)或\(m=-13\)。故\(m+n=3\)或\(-17\)，故选：AB．}
\end{ansblock}

\liB{变式\textbf{10}}{中档(平行线带域)}{}{在同一平面内，已知直线\(l_1\)：\(2x+3y=1\)和直线\(l_2\)：\(4x+6y-9=0\)，若直线\(l\)到直线\(l_1\)的距离与到直线\(l_2\)的距离之比为\(2:1\)，则直线\(l\)的方程为\kongbai{}．}
\xiexwei{9mm}

{\ansblockgrayfalse
\begin{ansblock}[2章-练-课时06B-T11]
% ans:2章-练-课时06B-T11
\ansitem{10}{2x+3y−8=0 或 6x+9y−10=0}
\ansnote{详解}{\(l_1\)化为\(4x+6y-2=0\)。\(l\)与两线均平行，设\(l\)：\(4x+6y+C=0\)（\(C\neq -2\)且\(C\neq -9\)）。由\(|C+2|/\sqrt{52}=2\times|C+9|/\sqrt{52}\)，即\(|C+2|=2|C+9|\)：\(C+2=2(C+9)\)得\(C=-16\)；\(C+2=-2(C+9)\)得\(C=-20/3\)。故\(l\)：\(4x+6y-16=0\)即\(2x+3y-8=0\)（带域外侧），或\(4x+6y-20/3=0\)即\(6x+9y-10=0\)（带域内侧，此时到\(l_1\)、\(l_2\)的距离分别为\(14/3\)与\(7/3\)，比值确为\(2:1\)）。两解均保留．}
\end{ansblock}}

\liB{变式\textbf{11}}{中档(距离几何意义·平行线间)}{}{已知\(m\)，\(n\)，\(a\)，\(b\in\mathbf{R}\)，且满足\(3m+4n=6\)，\(3a+4b=1\)，则\(\sqrt{(m-a)^{2}+(n-b)^{2}}\)的最小值为\kongbai{}．}
\xiexwei{7mm}

\begin{ansblock}[2章-练-课时06B-T14]
% ans:2章-练-课时06B-T14
\ansitem{11}{1}
\ansnote{详解}{设\(A(m,n)\)，\(B(a,b)\)，则\(A\)在直线\(3x+4y-6=0\)上，\(B\)在直线\(3x+4y-1=0\)上，所求即\(|AB|\)。两线平行，当\(AB\)沿公垂方向时最短，最小值为平行线间的距离：\(d=|-6-(-1)|/\sqrt{3^{2}+4^{2}}=5/5=1\)．}
\end{ansblock}

\xiaojie{①识别：平行线间距离、动平行线范围、平行线系成比、折线和最值，判归本题型；核心是“化同系数—套公式—看几何意义”三步.}
\bindp {\kaishu ②操作：系数不同先化同；范围问题抓两端——最大是公垂（或\(PQ\)长）、下界看是否可取0；平行线系设\(Ax+By+m=0\)后解\(m\).}
\bindp {\kaishu ③收束：垂线段转化后用“三角形不等式＋共线取等”，取等点必须回代验证确在约束直线上，防假最值.}

% ---------- 探究点三 对称·反射与拓展 ----------
\tjdnr{三}{对称·反射与拓展}{例\textbf{12}}{简单(光线反射)}{光线从点\(A(-3,5)\)射到\(x\)轴上，经反射以后经过点\(B(2,10)\)，则光线从\(A\)到\(B\)经过的路程为\nobreak (\kongwei)}

\bindopt {\fontsize{10.5pt}{\qpTJgLeadLiti}\selectfont \makebox[\dimexpr0.25\linewidth-0.5em\relax][l]{A．\(5\sqrt{2}\)}\makebox[\dimexpr0.25\linewidth-0.5em\relax][l]{B．\(2\sqrt{5}\)}\makebox[\dimexpr0.25\linewidth-0.5em\relax][l]{C．\(5\sqrt{10}\)}\makebox[\dimexpr0.25\linewidth-0.5em\relax][l]{D．\(10\sqrt{5}\)}\par}

\begin{ansblock}[2章-练-课时06B-T2]
% ans:2章-练-课时06B-T2
\ansitem{12}{C}
\ansnote{详解}{\(A\)关于反射面\(x\)轴的对称点为\(A'(-3,-5)\)。由反射角等于入射角，光路\(A\to M\to B\)的总长等于直线段\(|A'B|\)的最短值\(|A'B|=\sqrt{(2+3)^{2}+(10+5)^{2}}=\sqrt{250}=5\sqrt{10}\)，故选：C．}
\end{ansblock}

\liB{变式\textbf{13}}{中档(对称求周长最值)}{}{设定点\(A(3,1)\)，\(B\)是\(x\)轴上的动点，\(C\)是直线\(y=x\)上的动点，则\(\triangle ABC\)周长的最小值是\nobreak (\kongwei)}

\bindopt {\fontsize{10.5pt}{\qpTJgLeadLiti}\selectfont \makebox[\dimexpr0.25\linewidth-0.5em\relax][l]{A．\(\sqrt{5}\)}\makebox[\dimexpr0.25\linewidth-0.5em\relax][l]{B．\(2\sqrt{5}\)}\makebox[\dimexpr0.25\linewidth-0.5em\relax][l]{C．\(3\sqrt{5}\)}\makebox[\dimexpr0.25\linewidth-0.5em\relax][l]{D．\(\sqrt{10}\)}\par}

{\ansblockgrayfalse
\begin{ansblock}[2章-练-课时06B-T4]
% ans:2章-练-课时06B-T4
\ansitem{13}{B}
\ansnote{详解}{\(A(3,1)\)关于\(y=x\)的对称点\(A_1(1,3)\)（对称式\((x,y)\to(y,x)\)），关于\(x\)轴的对称点\(A_2(3,-1)\)。则\(|AC|=|A_1C|\)，\(|AB|=|A_2B|\)，周长\(=|A_1C|+|CB|+|BA_2|\geqslant|A_1A_2|\)。取等检验：直线\(A_1A_2\)为\(y=-2x+5\)，与\(y=x\)交于\((5/3,\ 5/3)\)、与\(x\)轴交于\((5/2,\ 0)\)，均落在允许动点范围内，等号可取。最小值\(=|A_1A_2|=\sqrt{(3-1)^{2}+(-1-3)^{2}}=\sqrt{20}=2\sqrt{5}\)，故选：B．}
\end{ansblock}}

\liB{变式\textbf{14}}{中档(含参对称点求\(a\)、\(b\))}{}{若点\(A(a+2,\ b+2)\)与\(B(b-a,\ -b)\)关于直线\(4x+3y-11=0\)对称，则实数\(a\)，\(b\)的值分别为\nobreak (\kongwei)}

\bindopt {\fontsize{10.5pt}{\qpTJgLeadLiti}\selectfont \makebox[\dimexpr0.25\linewidth-0.5em\relax][l]{A．\(-1\)，\(2\)}\makebox[\dimexpr0.25\linewidth-0.5em\relax][l]{B．\(4\)，\(-2\)}\makebox[\dimexpr0.25\linewidth-0.5em\relax][l]{C．\(2\)，\(4\)}\makebox[\dimexpr0.25\linewidth-0.5em\relax][l]{D．\(4\)，\(2\)}\par}

\begin{ansblock}[2章-练-课时06B-T7]
% ans:2章-练-课时06B-T7
\ansitem{14}{D}
\ansnote{详解}{对称\(\iff AB\perp l\)且\(AB\)的中点在\(l\)上。\(k_{AB}=(-b-(b+2))/((b-a)-(a+2))=(-2b-2)/(b-2a-2)=3/4\)，整理得\(-8b-8=3b-6a-6\)，即\(6a-11b-2=0\)①。中点坐标\(((b+2)/2,\ 1)\)，代入\(4x+3y-11=0\)：\(2(b+2)+3-11=0\)，得\(b=2\)；代回①：\(6a-22-2=0\)，\(a=4\)。故\(a=4\)，\(b=2\)，故选：D．}
\end{ansblock}

\tjdnr{三}{对称·反射与拓展}{例\textbf{15}}{中档(角平分线等距·对称)}{已知\(\triangle ABC\)的内角平分线\(CD\)的方程为\(2x+y-1=0\)，两个顶点为\(A(1,2)\)，\(B(-1,-1)\)。}

\bindp (1) 求点\(A\)到直线\(CD\)的距离；

\bindp (2) 求点\(C\)的坐标。
\xiexwei{18mm}

{\ansblockgrayfalse
\begin{ansblock}[2章-练-课时06B-T15]
% ans:2章-练-课时06B-T15
\ansitem{15}{（1）3√5/5；（2）C(−13/5, 31/5)。}
\ansnote{详解}{（1）\(d=|2\times 1+2-1|/\sqrt{2^{2}+1^{2}}=3/\sqrt{5}=3\sqrt{5}/5\)。\par\noindent （2）内角平分线性质：\(A\)关于\(CD\)的对称点\(A'\)落在边\(BC\)所在直线上。设\(A'(x,y)\)：中点\(((x+1)/2,\ (y+2)/2)\)在\(CD\)上，得\((x+1)+(y+2)/2-1=0\)，即\(2x+y+2=0\)；\(AA'\perp CD\)，\(k_{CD}=-2\)，故\(k_{AA'}=1/2\)，得\((y-2)/(x-1)=1/2\)，即\(x-2y+3=0\)。\par\noindent 联立解得\(x=-7/5\)，\(y=4/5\)，即\(A'(-7/5,\ 4/5)\)。直线\(BC\)过\(B(-1,-1)\)与\(A'\)，斜率\(=(4/5+1)/(-7/5+1)=(9/5)/(-2/5)=-9/2\)，方程\(y+1=-(9/2)(x+1)\)，即\(9x+2y+11=0\)。\par\noindent \(C\)为\(BC\)与\(CD\)的交点：把\(y=1-2x\)代入，\(9x+2(1-2x)+11=0\)，\(5x+13=0\)，\(x=-13/5\)，\(y=1-2\times(-13/5)=31/5\)。故\(C(-13/5,\ 31/5)\)．}
\end{ansblock}}

\liB{变式\textbf{16}}{中档(距离方程组探究)}{}{已知三条直线\(l_1\)：\(2x-y+3=0\)，\(l_2\)：\(-4x+2y+1=0\)和\(l_3\)：\(x+y-1=0\)。能否找到一点\(P\)，使得点\(P\)同时满足下列三个条件：}

\bindp (1) \(P\)是第一象限的点；

\bindp (2) \(P\)点到\(l_1\)的距离是\(P\)点到\(l_2\)的距离的\(1/2\)；

\bindp (3) \(P\)点到\(l_1\)的距离与\(P\)点到\(l_3\)的距离之比是\(\sqrt{2}:\sqrt{5}\)？

\bindp 若能，求出\(P\)点的坐标；若不能，说明理由。
\xiexwei{18mm}

{\ansblockgrayfalse
\begin{ansblock}[2章-练-课时06B-T16]
% ans:2章-练-课时06B-T16
\ansitem{16}{能，P(1/9, 37/18)。}
\ansnote{详解}{设\(P(x_0,\ y_0)\)，记\(u=2x_0-y_0\)。\par\noindent 条件（3）：\(|u+3|/\sqrt{5}=(\sqrt{2}/\sqrt{5})\cdot|x_0+y_0-1|/\sqrt{2}\)，即\(|u+3|=|x_0+y_0-1|\)，平方去绝对值：\((u+3)^{2}-(x_0+y_0-1)^{2}=(3x_0+2)(x_0-2y_0+4)=0\)。\(x_0=-2/3<0\)与第一象限矛盾，舍；故\(x_0-2y_0+4=0\)①。\par\noindent 条件（2）：\(d_1=|u+3|/\sqrt{5}\)，\(d_2=|-2u+1|/(2\sqrt{5})\)，由\(d_1=(1/2)d_2\)得\(4|u+3|=|1-2u|\)，平方：\((6u+11)(2u+13)=0\)，\(u=-11/6\)或\(u=-13/2\)，即\(2x_0-y_0+11/6=0\)②或\(2x_0-y_0+13/2=0\)③。\par\noindent 联立逐一检验：②与①：\(y_0=2x_0+11/6\)，代入①：\(x_0-4x_0-11/3+4=0\)，\(x_0=1/9>0\)，\(y_0=2/9+11/6=37/18>0\)，在第一象限；③与①：\(y_0=2x_0+13/2\)，代入①：\(x_0-4x_0-13+4=0\)，\(x_0=-3<0\)，舍。\par\noindent 终检\(P(1/9,\ 37/18)\)：\(d_1=(21/18)/\sqrt{5}=7/(6\sqrt{5})\)，\(d_2=(14/3)/(2\sqrt{5})=7/(3\sqrt{5})\)，\(d_1=(1/2)d_2\)成立；\(d_3=(21/18)/\sqrt{2}=7/(6\sqrt{2})\)，\(d_1:d_3=\sqrt{2}:\sqrt{5}\)成立。故存在满足三条件的点\(P(1/9,\ 37/18)\)．}
\end{ansblock}}

\xiaojie{①识别：角平分线（等距）、光线反射、对称点与对称最短，判归本题型；对称的两条约束（垂直＋中点在轴上）是通用工具.}
\bindp {\kaishu ②操作：反射／对称最短先作对称点，把折线拉直成两定点连线段；角平分线用“角顶两侧等距”或对称落边转化.}
\bindp {\kaishu ③收束：多条件探究把每个条件都化成含参方程（组），逐支求解、逐支检验（象限、重合、退化），最后终检回代各条件.}

% ============================================================
% 课堂评价 G5（恒5＝3单选+2填空＝导槽 G1~G5；ketang 新制顺排可跨栏，禁 \vbox 装栏）
% 印面号 17—21；多选门：本片正文多选 1（T10）≤2 合规
% ============================================================
\begin{ketang}{本组共5题（第17—21题），17—19为单选，20—21为填空．}

\jiancestem{{\fontsize{11.4pt}{14pt}\selectfont\heihao {\numboldjian 17}．}\hspace{4.1mm}\tieside{中档(距离公式直用求参)}\hspace{2.2mm}已知点\(M(1,4)\)到直线\(l\)：\(mx+y-1=0\)的距离为3，则实数\(m\)等于\nobreak(\kongwei)}

\bindopt {\fontsize{10.5pt}{\qpTJgLeadPingjia}\selectfont \makebox[\dimexpr0.5\linewidth-1em\relax][l]{A．\(0\)}\makebox[\dimexpr0.5\linewidth-1em\relax][l]{B．\(3/4\)}\par}

\bindopt {\fontsize{10.5pt}{\qpTJgLeadPingjia}\selectfont \makebox[\dimexpr0.5\linewidth-1em\relax][l]{C．\(3\)}\makebox[\dimexpr0.5\linewidth-1em\relax][l]{D．\(0\)或\(3/4\)}\par}

\begin{ansblock}[2章-导-课时06B-G1]
% ans:2章-导-课时06B-G1
\ansitem{17}{D}
\ansnote{详解}{\(d=|m\cdot 1+4-1|/\sqrt{m^{2}+1}=|m+3|/\sqrt{m^{2}+1}=3\)，两边非负可平方：\(m^{2}+6m+9=9m^{2}+9\)，即\(8m^{2}-6m=0\)，\(m=0\)或\(m=3/4\)。检验：\(m=0\)时直线为\(y-1=0\)，\(d=|4-1|=3\)成立；\(m=3/4\)时\(d=(15/4)/(5/4)=3\)成立。两根均保留，故选：D．}
\end{ansblock}

\jiancestem{{\fontsize{11.4pt}{14pt}\selectfont\heihao {\numboldjian 18}．}\hspace{4.1mm}\tieside{简单(公式直用)}\hspace{2.2mm}原点\(O(0,0)\)到直线\(3x+4y-5=0\)的距离为\nobreak(\kongwei)}

\bindopt {\fontsize{10.5pt}{\qpTJgLeadPingjia}\selectfont \makebox[\dimexpr0.25\linewidth-0.5em\relax][l]{A．\(1\)}\makebox[\dimexpr0.25\linewidth-0.5em\relax][l]{B．\(\sqrt{5}\)}\makebox[\dimexpr0.25\linewidth-0.5em\relax][l]{C．\(5\)}\makebox[\dimexpr0.25\linewidth-0.5em\relax][l]{D．\(1/5\)}\par}

\begin{ansblock}[2章-导-课时06B-G2]
% ans:2章-导-课时06B-G2
\ansitem{18}{A}
\ansnote{详解}{\(d=|3\times 0+4\times 0-5|/\sqrt{3^{2}+4^{2}}=5/5=1\)，故选：A．}
\end{ansblock}

\jiancestem{{\fontsize{11.4pt}{14pt}\selectfont\heihao {\numboldjian 19}．}\hspace{4.1mm}\tieside{中档(平行判定与平行线距离)}\hspace{2.2mm}若直线\(l_1\)：\(x+ay+6=0\)与\(l_2\)：\((a-2)x+3y+2a=0\)平行，则\(l_1\)与\(l_2\)间的距离为\nobreak(\kongwei)}

\bindopt {\fontsize{10.5pt}{\qpTJgLeadPingjia}\selectfont \makebox[\dimexpr0.25\linewidth-0.5em\relax][l]{A．\(\sqrt{2}\)}\makebox[\dimexpr0.25\linewidth-0.5em\relax][l]{B．\(8\sqrt{2}/3\)}\makebox[\dimexpr0.25\linewidth-0.5em\relax][l]{C．\(\sqrt{3}\)}\makebox[\dimexpr0.25\linewidth-0.5em\relax][l]{D．\(8\sqrt{3}/3\)}\par}

\begin{ansblock}[2章-导-课时06B-G3]
% ans:2章-导-课时06B-G3
\ansitem{19}{B}
\ansnote{详解}{平行\(\iff 1\times 3-a(a-2)=0\)且系数不成比例（排除重合），由\(a^{2}-2a-3=0\)得\(a=3\)或\(a=-1\)；\(a=3\)时两方程同为\(x+3y+6=0\)（重合），舍，故\(a=-1\)。此时\(l_1\)：\(x-y+6=0\)，\(l_2\)：\(-3x+3y-2=0\)即\(x-y+2/3=0\)，\(d=|6-2/3|/\sqrt{1^{2}+(-1)^{2}}=(16/3)/\sqrt{2}=8\sqrt{2}/3\)，故选：B．}
\end{ansblock}

\jiancestem{{\fontsize{11.4pt}{14pt}\selectfont\heihao {\numboldjian 20}．}\hspace{4.1mm}\tieside{简单(平行线距离·系数化)}\hspace{2.2mm}求平行线\(2x-3y+5=0\)与\(2x-3y-8=0\)之间的距离：\kongbai{}．}
\xiexwei{7mm}

\begin{ansblock}[2章-导-课时06B-G4]
% ans:2章-导-课时06B-G4
\ansitem{20}{√13}
\ansnote{详解}{两线\(x\)、\(y\)系数已相同，直接套平行线间的距离公式：\(d=|5-(-8)|/\sqrt{2^{2}+(-3)^{2}}=13/\sqrt{13}=\sqrt{13}\)．}
\end{ansblock}

\jiancestem{{\fontsize{11.4pt}{14pt}\selectfont\heihao {\numboldjian 21}．}\hspace{4.1mm}\tieside{简单(距离求参·先化一般式)}\hspace{2.2mm}已知点\(B(m,6)\)到直线\(y=3x+6\)的距离为3，则实数\(m\)的值为\kongbai{}．}
\xiexwei{7mm}

\begin{ansblock}[2章-导-课时06B-G5]
% ans:2章-导-课时06B-G5
\ansitem{21}{m＝±√10}
\ansnote{详解}{\(y=3x+6\)化为一般式\(3x-y+6=0\)，\(d=|3m-6+6|/\sqrt{3^{2}+(-1)^{2}}=3|m|/\sqrt{10}=3\)，得\(|m|=\sqrt{10}\)，即\(m=\pm\sqrt{10}\)．}
\end{ansblock}

\end{ketang}

% ---- 尾块（\tailfill 置 \end{multicols} 前；无参＝内容尾随框，件侧不擅自定高） ----
\tailfill

\end{multicols}
\end{document}
